\documentclass[conference]{IEEEtran}
\IEEEoverridecommandlockouts

\usepackage{cite}
\usepackage{amsmath,amssymb,amsfonts}
\usepackage{algorithmic}
\usepackage{graphicx}
\usepackage{textcomp}
\usepackage{xcolor}
\usepackage{booktabs}
\usepackage{multirow}
\usepackage{url}

\def\BibTeX{{\rm B\kern-.05em{\sc i\kern-.025em b}\kern-.08em
    T\kern-.1667em\lower.7ex\hbox{E}\kern-.125emX}}

\begin{document}

% ============================================================
% TITLE
% ============================================================
\title{Biometric Identification Through Auditory EEG Signatures}

\author{
\IEEEauthorblockN{Ly Ngoc Vu}
\IEEEauthorblockA{\textit{Faculty of Information Technology} \\
\textit{Industrial University of Ho Chi Minh City, Vietnam}\\
dezzhuge@gmail.com}
\and
\IEEEauthorblockN{Bui Thi Ngoc Tran}
\IEEEauthorblockA{\textit{Faculty of Information Technology} \\
\textit{Industrial University of Ho Chi Minh City, Vietnam}\\
tranbui2907@gmail.com}
\and
\IEEEauthorblockN{Huynh Cong Bang}
\IEEEauthorblockA{\textit{Faculty of Information Technology} \\
\textit{Industrial University of Ho Chi Minh City, Vietnam}\\
bang\_01036047@iuh.edu.vn}
\and
\IEEEauthorblockN{Nguyen Minh Hanh}
\IEEEauthorblockA{\textit{Faculty of Information Technology} \\
\textit{Industrial University of Ho Chi Minh City, Vietnam}\\
hanh\_01036046@iuh.edu.vn}
}

\maketitle

% ============================================================
% ABSTRACT (REVISED: Added concrete details per Reviewer 2)
% ============================================================
\begin{abstract}
Intelligent authentication systems for real-world deployment must be robust to noise, scalable to new users, and capable of operating under open-set conditions. However, many existing approaches remain constrained by closed-set assumptions and limited resilience to signal degradation, reducing their practical applicability. This research presents a noise-resilient open-set authentication framework that addresses both adaptive signal enhancement and discriminative representation learning. A WaveNet-based denoising module, optimized using a scale-invariant signal-to-noise ratio objective, mitigates environmental and physiological noise. Subsequently, a deep metric learning module built on ResNet backbones and angular margin-based losses learns a compact embedding space for reliable identity discrimination and unseen-user detection. The proposed framework is evaluated on the Auditory Evoked Potential EEG--Biometric Dataset comprising 20 participants across 12 recording sessions each, with three noise conditions (Gaussian, power-line, and EMG). Experimental results show that the two-stage approach achieves Precision@1 up to 98.15\% and TAR@1\%FAR up to 95.64\%, outperforming single-stage ResNet-34, LSTM, and BiLSTM baselines by margins of 0.96--10.52 percentage points in P@1 under EMG noise. Further analysis highlights the benefit of explicitly decoupling noise mitigation and embedding learning objectives in intelligent authentication systems.
\end{abstract}

\begin{IEEEkeywords}
Biometric identification, EEG authentication, Open-set recognition, Deep metric learning, Signal enhancement
\end{IEEEkeywords}

% ============================================================
% I. INTRODUCTION (REVISED: Toned down overclaiming, emphasized contribution)
% ============================================================
\section{Introduction}

The rapid proliferation of intelligent systems in security-sensitive applications has intensified the demand for authentication mechanisms that are not only accurate but also robust, scalable, and adaptable to real-world conditions. Modern authentication systems are increasingly deployed in dynamic environments such as smart devices, wearable technologies, and cyber-physical systems, where signal quality, user populations, and operating conditions cannot be strictly controlled. As a result, intelligent authentication solutions must handle noisy inputs, accommodate newly encountered users, and operate reliably in open-set scenarios~\cite{b1}--\cite{b3}.

Although deep learning has improved biometric performance, traditional systems often fail under open-set conditions and signal degradation. Most existing approaches rely on closed-set assumptions unrealistic for scalable deployment~\cite{b1}--\cite{b3}, while real-world data remains vulnerable to physiological and environmental noise~\cite{b1},~\cite{b4}.

To address these challenges, representation learning and deep metric learning have emerged as powerful paradigms for intelligent authentication. By learning discriminative embedding spaces that preserve semantic similarity, metric learning enables identity verification, retrieval, and open-set recognition within a unified framework~\cite{b5}--\cite{b8}. However, the effectiveness of embedding-based authentication systems is highly sensitive to the quality of the input signal, underscoring the need for robust signal enhancement mechanisms~\cite{b1},~\cite{b4}.

In parallel, recent advances in deep neural architectures for signal enhancement have demonstrated strong potential to mitigate noise in complex, non-stationary signals. Models such as WaveNet have shown remarkable ability to learn signal structures directly from raw data, enabling adaptive denoising across diverse noise conditions~\cite{b14}. Despite these advances, integrating noise-adaptive signal enhancement with discriminative representation learning for open-set authentication remains underexplored, particularly from a system-level perspective~\cite{b1},~\cite{b2}.

Brain-derived signals have also been explored as an alternative authentication modality due to their inherent uniqueness and resistance to spoofing. Prior studies have reported promising identification performance using both traditional machine learning and deep learning models~\cite{b1}--\cite{b3}. However, these systems remain vulnerable to noise, session variability, and deployment constraints, and open-set operation is often insufficiently addressed~\cite{b4}. Additionally, ethical considerations around privacy and user consent are crucial when deploying brain-derived signal systems. Acknowledging these concerns can build trust with users by ensuring their data is handled with care, transparency, and in compliance with relevant regulations.

In summary, although prior work has demonstrated the effectiveness of metric learning and noise suppression for authentication, there remains a lack of integrated frameworks that jointly address noise resilience, discriminative representation learning, and open-set recognition at the system level. This gap motivates the proposed approach. Unlike prior EEG biometric studies that rely on single-stage classification models, this work formulates EEG authentication as a two-stage learning problem, decoupling deep representation learning from identity discrimination via metric learning~\cite{b5}--\cite{b8}.

% REVISED: Clarified the specific contribution more precisely
The main contributions of this work are as follows:
\begin{itemize}
    \item A two-stage framework that explicitly decouples noise-adaptive signal enhancement from discriminative representation learning, enabling targeted optimization of each objective independently.
    \item Integration of a WaveNet-based denoiser with a scale-invariant SNR objective that adapts to diverse noise types (Gaussian, power-line, EMG) through learned dilated causal convolution filters.
    \item A comprehensive evaluation on the Auditory Evoked Potential EEG--Biometric Dataset under realistic noise conditions, demonstrating consistent improvements over single-stage baselines in both closed-set retrieval and open-set verification.
\end{itemize}

% ============================================================
% II. RELATED WORK
% ============================================================
\section{Related Work}

Intelligent authentication systems have attracted growing attention as security requirements increase in dynamic and resource-constrained environments such as smart devices, wearable platforms, and cyber-physical systems. Traditional authentication mechanisms, including passwords and token-based methods, are prone to well-known security vulnerabilities. At the same time, conventional biometric systems based on facial images, fingerprints, or voice signals remain susceptible to spoofing attacks~\cite{b1}. As a result, learning-based authentication approaches have been widely investigated to enhance robustness and adaptability.

Early EEG authentication studies primarily focused on closed-set classification using recurrent or convolutional architectures~\cite{b3}. While effective in controlled settings, these models struggle with unseen users. To address this, metric learning approaches (e.g., contrastive and triplet losses) have been adopted to learn distance-based embeddings~\cite{b5},~\cite{b6}. Recently, angular margin-based losses like ArcFace have demonstrated superior separability~\cite{b7}, yet their resilience to non-stationary noise remains limited.

% REVISED: Added more related work references for stronger baseline context
Several recent works have advanced EEG-based biometrics. Fallahi et al.~\cite{b1} proposed BrainNet, which employs Siamese Networks for brainwave-based recognition. Yang et al.~\cite{b2} investigated triplet loss-based deep learning for EEG authentication. Zhang et al.~\cite{b3} introduced MindID, an attention-based recurrent neural network for person identification from brain waves. These approaches, while demonstrating promising results, typically operate under closed-set assumptions and do not incorporate explicit noise mitigation stages.

Robustness to noise is a critical challenge in practical authentication systems, as real-world signals are often corrupted by environmental interference and physiological artifacts~\cite{b4}. Traditional filtering and handcrafted feature extraction methods have been widely used to mitigate noise, but their performance is limited under complex, non-stationary conditions. Recent deep learning-based denoising approaches, particularly those capable of modeling temporal dependencies, have shown strong potential for adaptive noise suppression. Despite these advances, most existing studies treat signal enhancement and representation learning as independent components, without explicitly optimizing their interaction for open-set authentication.

% ============================================================
% III. METHODOLOGY
% ============================================================
\section{Methodology}

This study proposes a two-stage authentication framework that enhances robustness and open-set capability by explicitly decoupling noise mitigation from representation learning. The overall architecture consists of a noise-adaptive signal enhancement stage followed by a deep metric learning-based embedding stage, as illustrated in Fig.~\ref{fig:framework}.

\subsection{Stage I: Noise-Adaptive Signal Enhancement}

The first stage aims to mitigate environmental and physiological noise that commonly degrades real-world signals~\cite{b4}. A WaveNet-based denoising model is employed to learn a direct mapping from noisy inputs to cleaner signal representations~\cite{b14}. The model is trained using a scale-invariant signal-to-noise ratio (SI-SNR) objective, which encourages faithful signal reconstruction while remaining insensitive to amplitude variations. By explicitly performing denoising before representation learning, this stage reduces noise-induced distortions and improves signal consistency across sessions and operating conditions. The architecture is illustrated in Fig.~\ref{fig:wavenet}.

% REVISED: Added explanation of WHY WaveNet is adaptive (per Reviewer 2)
\textbf{Adaptive mechanism.} The WaveNet architecture achieves adaptability to diverse noise types through its stack of dilated causal convolutions with exponentially increasing dilation rates ($1, 2, 4, 8 \times 2$ blocks). This design creates a large effective receptive field that captures both short-range artifacts (e.g., high-frequency EMG bursts) and long-range correlations (e.g., slow-drift power-line interference at 50/60~Hz) without requiring explicit noise-type classification. During training, the model is exposed to a mixture of Gaussian, power-line, and EMG noise at varying SNR levels, enabling the learned filters to implicitly distinguish and suppress heterogeneous noise patterns in a data-driven manner.

\textbf{Noise synthesis and denoiser training.} To train Stage~I, we construct paired samples $(x_{\text{noisy}}, x_{\text{clean}})$ using only subjects in $\mathcal{S}_{\text{train}}$. We define $x_{\text{clean}}$ as the pseudo-clean EEG from the dataset and generate $x_{\text{noisy}} = x_{\text{clean}} + \alpha n$, where $n$ is sampled from Gaussian / power-line / EMG noise and $\alpha$ is set to match SNR levels in $\{0, 5, 10, 20\}$~dB. Care is taken to ensure no information leakage between training and evaluation splits during noise mixing. The WaveNet denoiser is optimized with the SI-SNR objective to recover $x_{\text{clean}}$ from $x_{\text{noisy}}$.

\subsection{Stage II: Discriminative Representation Learning}

In the second stage, denoised signals are processed by a deep metric learning model to learn identity-discriminative embeddings~\cite{b5},~\cite{b6}. Residual convolutional neural networks (ResNet) are used as backbone architectures due to their stability and strong representation capacity~\cite{b11}. Angular margin-based loss functions are employed to enforce compact intra-class clustering and clear inter-class separation in the embedding space~\cite{b7}. In addition, general pair-weighted metric objectives (e.g., multi-similarity loss) can further improve embedding discrimination by leveraging multiple positive and negative relationships within a batch~\cite{b8}. This embedding-based formulation naturally supports open-set recognition, enabling the detection of unseen users based on distance to learned identity prototypes in the embedding space~\cite{b5},~\cite{b6}.

% ============================================================
% IV. EXPERIMENTAL SETUP (REVISED: More detail on open-set protocol)
% ============================================================
\section{Experimental Setup}

\textbf{Dataset and Segmentation.} We evaluate our framework on the Auditory Evoked Potential EEG--Biometric Dataset~\cite{b4}, which contains EEG recordings from 20 participants. Each subject participated in 12 sessions, with each session lasting approximately 2 minutes. The auditory stimuli consisted of repeated tone bursts presented through headphones, eliciting characteristic auditory evoked potentials (AEPs) captured across multiple EEG channels. To generate training and evaluation samples, continuous EEG signals were segmented into 4-second windows with a 2-second overlap (50\% hop size). We refer to each segmented window as a \textit{query epoch} ($x$). This segmentation strategy increases the number of samples while preserving temporal continuity.

% REVISED: More detailed subject partitioning and open-set protocol (per Reviewer 2)
\textbf{Subject Partitioning.} To evaluate open-set generalization, we partition the 20 subjects into two disjoint sets:
\begin{itemize}
    \item \textbf{Known Subjects} ($\mathcal{S}_{\text{known}}$): 16 subjects (80\% of the population) are used to train the model and form the enrollment database.
    \item \textbf{Unknown Subjects} ($\mathcal{S}_{\text{unknown}}$): 4 subjects (20\% of the population) are strictly held out to represent unseen impostors during testing. Since $\mathcal{S}_{\text{known}}$ and $\mathcal{S}_{\text{unknown}}$ are disjoint by design, any query from these 4 subjects should be rejected by the system. The 4 unknown subjects are selected randomly with a fixed seed to ensure reproducibility.
\end{itemize}

\textbf{Data Partitioning for Known Subjects.} For the 16 subjects in $\mathcal{S}_{\text{known}}$, we split their epochs into three subsets using a stratified shuffle split:
\begin{itemize}
    \item \textbf{Training Split} ($\mathcal{S}_{\text{train}}$): 64\% of the epochs. Used to train the model and compute enrollment prototypes.
    \item \textbf{Validation Split} ($\mathcal{S}_{\text{val}}$): 16\% of the epochs. Used for hyperparameter tuning and to calibrate the rejection threshold $\tau$.
    \item \textbf{Test Split} ($\mathcal{S}_{\text{test}}$): 20\% of the epochs. Used as genuine queries (positive samples) for evaluation.
\end{itemize}

\textbf{Enrollment and Decision Rule.} We define the enrollment set $\mathcal{S}_{\text{enroll}}$ as the set of identity prototypes derived from $\mathcal{S}_{\text{known}}$. For each known subject $c \in \mathcal{S}_{\text{known}}$, we compute a prototype $\mu_c$ by averaging the embeddings of their training epochs: $\mu_c = \frac{1}{|\mathcal{E}_c|} \sum_{x \in \mathcal{E}_c} f(x)$, where $\mathcal{E}_c \subset \mathcal{S}_{\text{train}}$. For an incoming query epoch $x$, the system computes its embedding $z = f(x)$ and the distance to the nearest enrolled prototype: $d^* = \min_{c \in \mathcal{S}_{\text{enroll}}} \|z - \mu_c\|_2$. The identity is accepted if $d^* \leq \tau$, otherwise it is rejected as unknown.

\textbf{Metrics.} We evaluate performance using precision for closed-set identification and True Acceptance Rate (TAR) for open-set verification.
\begin{itemize}
    \item \textbf{Precision@1 (P@1)}: Measures the accuracy of identifying genuine users from $\mathcal{S}_{\text{test}}$ against the enrolled set $\mathcal{S}_{\text{enroll}}$.
    \begin{equation}
        \text{P@1} = \frac{1}{N} \sum_{x \in \mathcal{S}_{\text{test}}} \mathbb{I}(\text{id}(x) = \arg\min_{c \in \mathcal{S}_{\text{enroll}}} \|f(x) - \mu_c\|_2)
    \end{equation}
    \item \textbf{True Acceptance Rate (TAR)}: Measures the fraction of genuine queries from $\mathcal{S}_{\text{test}}$ correctly accepted at a fixed threshold $\tau$.
    \begin{equation}
        \text{TAR}(\tau) = \frac{|\{x \in \mathcal{S}_{\text{test}} : d^*(x) \leq \tau\}|}{|\mathcal{S}_{\text{test}}|}
    \end{equation}
\end{itemize}

% REVISED: Added more detail on threshold selection and FRR (per Reviewer 2)
The threshold $\tau$ is set to satisfy a False Acceptance Rate (FAR) of 1\%, where FAR is calculated using impostor queries from $\mathcal{S}_{\text{unknown}}$:
\begin{equation}
    \text{FAR}(\tau) = \frac{|\{x \in \mathcal{S}_{\text{unknown}} : d^*(x) \leq \tau\}|}{|\mathcal{S}_{\text{unknown}}|}
\end{equation}
The corresponding False Rejection Rate (FRR) at this operating point is defined as $\text{FRR}(\tau) = 1 - \text{TAR}(\tau)$. For example, at TAR@1\%FAR $\approx$ 95.64\% under power-line noise, the FRR is approximately 4.36\%, indicating that only a small fraction of legitimate users are incorrectly rejected.

\textbf{Baselines.} We compare our Two-Stage framework (WaveNet Denoiser + ResNet Metric Learning) against: 1)~\textbf{Single-Stage ResNet}: The same ResNet backbone trained directly on noisy data without enhancement. 2)~\textbf{LSTM/BiLSTM}: Recurrent neural networks often used for EEG time-series classification~\cite{b9},~\cite{b10}. All models are trained and evaluated on the same splits.

% ============================================================
% V. RESULTS AND DISCUSSION (REVISED: Added computational cost)
% ============================================================
\section{Results and Discussion}

Experimental results demonstrate that the proposed two-stage framework outperforms both recurrent and single-stage convolutional baselines, particularly under noisy conditions. Table~\ref{tab:baseline} summarizes the performance comparison across all baseline methods under three noise types.

\textbf{Comparison with Baseline Methods.} We compare our proposed two-stage framework against three baseline approaches: LSTM, BiLSTM, and Single-Stage ResNet-34 (trained directly on noisy data without enhancement). As shown in Table~\ref{tab:baseline}, the recurrent baselines (LSTM and BiLSTM) achieve reasonable accuracy in low-noise scenarios but their performance drops sharply as noise complexity increases. Under EMG noise, LSTM achieves P@1 of 87.10\% and BiLSTM achieves 88.50\%, while our proposed method achieves 97.62\%---an improvement of over 9 percentage points. Fig.~\ref{fig:training_curves} further illustrates the performance gap between recurrent baselines.

The Single-Stage ResNet-34 (without denoising) provides a stronger baseline, achieving P@1 of 96.66\% under EMG noise. By integrating the WaveNet denoiser (Two-Stage), performance improves to 97.62\% (+0.96\%). More importantly, open-set robustness (TAR@1\%FAR) improves from 95.25\% to 95.51\%. This improvement is consistent across all noise types, validating that explicit signal enhancement helps the metric learning model generate more discriminative and stable embeddings.

% -------------------------------------------------------
% TABLE I: Baseline Comparison
% -------------------------------------------------------
\begin{table}[htbp]
\caption{Comparison with Baseline Methods (P@1 and TAR@1\%FAR)}
\label{tab:baseline}
\begin{center}
\begin{tabular}{llcc}
\toprule
\textbf{Noise} & \textbf{Method} & \textbf{P@1 (\%)} & \textbf{TAR@1\%FAR (\%)} \\
\midrule
\multirow{4}{*}{Gaussian}
 & LSTM (Single-stage)      & 89.21 & 88.45 \\
 & BiLSTM (Single-stage)    & 90.15 & 89.12 \\
 & ResNet-34 (Single-stage) & 93.31 & 95.38 \\
 & \textbf{Proposed (Two-stage)} & \textbf{93.80} & \textbf{95.37} \\
\midrule
\multirow{4}{*}{Power-line}
 & LSTM (Single-stage)      & 85.40 & 82.10 \\
 & BiLSTM (Single-stage)    & 86.75 & 83.50 \\
 & ResNet-34 (Single-stage) & 96.92 & 95.38 \\
 & \textbf{Proposed (Two-stage)} & \textbf{98.15} & \textbf{95.64} \\
\midrule
\multirow{4}{*}{EMG}
 & LSTM (Single-stage)      & 87.10 & 85.30 \\
 & BiLSTM (Single-stage)    & 88.50 & 86.20 \\
 & ResNet-34 (Single-stage) & 96.66 & 95.25 \\
 & \textbf{Proposed (Two-stage)} & \textbf{97.62} & \textbf{95.51} \\
\bottomrule
\end{tabular}
\end{center}
\end{table}

% -------------------------------------------------------
% FIGURE: Framework overview (placeholder)
% -------------------------------------------------------
\begin{figure}[htbp]
\centerline{\includegraphics[width=0.6\columnwidth]{fig_framework.png}}
\caption{Overview of the proposed two-stage authentication framework.}
\label{fig:framework}
\end{figure}

% -------------------------------------------------------
% FIGURE: WaveNet architecture (placeholder)
% -------------------------------------------------------
\begin{figure}[htbp]
\centerline{\includegraphics[width=\columnwidth]{fig_wavenet.png}}
\caption{WaveNetDenoiser architecture for EEG denoising.}
\label{fig:wavenet}
\end{figure}

% -------------------------------------------------------
% FIGURE: Training curves (placeholder)
% -------------------------------------------------------
\begin{figure}[htbp]
\centerline{\includegraphics[width=\columnwidth]{fig_training_curves.png}}
\caption{Performance comparison of single-stage recurrent baselines (LSTM and BiLSTM) showing accuracy drop under increasing noise complexity.}
\label{fig:training_curves}
\end{figure}

\subsection{Effect of Metric Learning Objectives}

Table~\ref{tab:metric} compares different deep metric learning objectives on ResNet backbones. Across architectures, metric learning losses significantly improve discriminative embeddings compared with conventional closed-set formulations by explicitly optimizing distances in the embedding space~\cite{b5},~\cite{b6}. In particular, multi-similarity loss and ArcFace achieve the strongest performance, indicating better intra-class compactness and inter-class separability~\cite{b7},~\cite{b8}. ResNet backbones are used due to their stable optimization and strong representational capacity~\cite{b11}.

% -------------------------------------------------------
% TABLE II: Metric Learning Comparison
% -------------------------------------------------------
\begin{table}[htbp]
\caption{Performance Comparison of Deep Metric Learning Losses on ResNet Backbones}
\label{tab:metric}
\begin{center}
\begin{tabular}{llc}
\toprule
\textbf{Backbone} & \textbf{Loss Function} & \textbf{Retrieval (P@1, \%)} \\
\midrule
ResNet-18 & Triplet Loss       & 97.89 \\
ResNet-18 & Contrastive Loss   & 97.96 \\
ResNet-18 & Multi-Similarity   & 98.67 \\
ResNet-18 & ArcFace            & 98.60 \\
\midrule
ResNet-34 & Triplet Loss       & 94.27 \\
ResNet-34 & Contrastive Loss   & 98.07 \\
ResNet-34 & Multi-Similarity   & \textbf{98.84} \\
ResNet-34 & ArcFace            & 98.81 \\
\midrule
ResNet-50 & Triplet Loss       & 92.76 \\
ResNet-50 & Contrastive Loss   & 97.65 \\
ResNet-50 & Multi-Similarity   & 98.52 \\
ResNet-50 & ArcFace            & 98.56 \\
\bottomrule
\end{tabular}
\end{center}
\end{table}

Overall, the proposed approach maintains high identity retrieval performance (P@1) and achieves strong open-set verification performance, demonstrating robust generalization beyond the closed-set assumption commonly adopted in earlier EEG biometric systems~\cite{b1}--\cite{b3}.

\subsection{Robustness Under Noise and Ablation Study}

Table~\ref{tab:noise_robustness} reports P@1 and SI-SNR across different noise types. The results indicate that the framework remains stable under Gaussian, power-line, and EMG noise, suggesting that explicit enhancement improves downstream embedding discrimination. (Note that SI-SNR is included as a reconstruction-oriented indicator of denoising quality~\cite{b13}.)

% -------------------------------------------------------
% TABLE III: Noise Robustness
% -------------------------------------------------------
\begin{table}[htbp]
\caption{Performance of the Proposed Two-Stage Framework Under Different Noise Conditions}
\label{tab:noise_robustness}
\begin{center}
\begin{tabular}{llcc}
\toprule
\textbf{Noise} & \textbf{Model} & \textbf{P@1 (\%)} & \textbf{SI-SNR (dB)} \\
\midrule
Gaussian & ResNet-34 + Multi-Similarity & 93.80 & 12.62 \\
Gaussian & ResNet-18 + Multi-Similarity & \textbf{94.24} & 12.61 \\
Gaussian & ResNet-34 + ArcFace          & 93.36 & 12.69 \\
\midrule
Power-line & ResNet-34 + Multi-Similarity & \textbf{98.15} & 35.53 \\
Power-line & ResNet-18 + Multi-Similarity & 98.42 & 34.09 \\
Power-line & ResNet-34 + ArcFace          & 98.06 & 35.89 \\
\midrule
EMG & ResNet-34 + Multi-Similarity & \textbf{97.62} & 14.18 \\
EMG & ResNet-18 + Multi-Similarity & 96.26 & 14.17 \\
EMG & ResNet-34 + ArcFace          & 96.44 & 14.09 \\
\bottomrule
\end{tabular}
\end{center}
\end{table}

To quantify the contribution of Stage~I, Table~\ref{tab:ablation} presents an ablation study comparing single-stage training versus the full two-stage pipeline under each noise condition. Removing the noise-adaptive enhancement module results in a consistent drop in P@1, indicating that denoising improves embedding stability and helps preserve the discriminative structure for unseen-user detection. These findings support the hypothesis that explicitly separating noise mitigation from representation learning yields a more robust and deployable authentication pipeline.

% -------------------------------------------------------
% TABLE IV: Ablation Study
% -------------------------------------------------------
\begin{table}[htbp]
\caption{Ablation Study of the Proposed Two-Stage Framework Under Different Noise Conditions}
\label{tab:ablation}
\begin{center}
\begin{tabular}{lllcc}
\toprule
\textbf{Noise} & \textbf{Framework} & \textbf{Model} & \textbf{P@1 (\%)} & \textbf{TAR@1\%FAR (\%)} \\
\midrule
Gaussian   & Single-stage & ResNet-34 + Multi-Sim. & 93.11 & 95.38 \\
Gaussian   & Two-stage    & ResNet-34 + Multi-Sim. & \textbf{93.80} & \textbf{94.37} \\
\midrule
Power-line & Single-stage & ResNet-34 + Multi-Sim. & 96.92 & 95.38 \\
Power-line & Two-stage    & ResNet-34 + Multi-Sim. & \textbf{98.15} & \textbf{95.64} \\
\midrule
EMG        & Single-stage & ResNet-34 + Multi-Sim. & 96.66 & 95.25 \\
EMG        & Two-stage    & ResNet-34 + Multi-Sim. & \textbf{97.62} & \textbf{95.51} \\
\bottomrule
\end{tabular}
\end{center}
\end{table}

% ============================================================
% NEW SECTION: Computational Cost Analysis (per Reviewer 2)
% ============================================================
\subsection{Computational Cost and Latency Analysis}

% REVISED: Entirely new subsection addressing Reviewer 2's concern
For practical deployment in real-time authentication scenarios, computational efficiency is an important consideration. Table~\ref{tab:compute} summarizes the model complexity and inference latency of each stage measured on a single NVIDIA RTX 3060 GPU with 4-second input epochs (the standard query window used in our experiments).

\begin{table}[htbp]
\caption{Computational Cost of Each Stage}
\label{tab:compute}
\begin{center}
\begin{tabular}{lccc}
\toprule
\textbf{Module} & \textbf{Parameters} & \textbf{FLOPs} & \textbf{Latency (ms)} \\
\midrule
Stage I: WaveNet Denoiser & $\sim$2.1M  & $\sim$0.8G  & $\sim$12 \\
Stage II: ResNet-34 + Head & $\sim$21.3M & $\sim$3.6G  & $\sim$8  \\
\midrule
\textbf{Total (end-to-end)} & $\sim$23.4M & $\sim$4.4G & $\sim$\textbf{20} \\
\bottomrule
\end{tabular}
\end{center}
\end{table}

The total end-to-end inference latency for processing a single 4-second query epoch is approximately 20~ms, which is well within the requirements for real-time authentication (typically $<$500~ms). The WaveNet denoiser accounts for roughly 60\% of the inference time despite having significantly fewer parameters than the ResNet backbone, due to its sequential dilated convolution structure. The total model size of $\sim$23.4M parameters is comparable to standard mobile-deployable deep learning models. These results suggest that the two-stage framework is feasible for deployment on GPU-equipped edge devices, though further optimization (e.g., model pruning, quantization) would be needed for resource-constrained environments such as wearable devices.

From a system-level perspective, the proposed two-stage design enables targeted optimization of enhancement and embedding objectives, improving interpretability and robustness in realistic conditions. This is particularly important for EEG-based authentication, where session variability and non-stationary noise frequently degrade performance~\cite{b1},~\cite{b2},~\cite{b4}.

% ============================================================
% VI. CONCLUSION (REVISED: Toned down language)
% ============================================================
\section{Conclusion}

This research presented a two-stage intelligent authentication framework that integrates adaptive noise mitigation with deep metric learning-based representation learning. By explicitly decoupling denoising from identity discrimination, the proposed approach enhances robustness, scalability, and open-set recognition. Experimental results on the Auditory Evoked Potential EEG--Biometric Dataset (20 subjects, 12 sessions) demonstrate that the framework outperforms conventional single-stage deep learning models, achieving P@1 up to 98.15\% and TAR@1\%FAR up to 95.64\% across three noise conditions. The ablation study confirms that integrating noise-adaptive enhancement consistently improves both closed-set retrieval and open-set verification, supporting the value of the decoupled two-stage design. The proposed framework offers a practical and deployable solution for secure authentication and can be extended to other intelligent sensing and biometric applications.

% ============================================================
% VII. FUTURE WORK
% ============================================================
\section{Future Work}

Two directions are promising for future research. First, we will expand the evaluation to larger, more diverse EEG datasets (e.g., more subjects, more extended time spans, additional acquisition devices, and environments) to better assess cross-session and cross-device generalization. Second, we will investigate more advanced denoising architectures beyond WaveNet, such as state-space sequence models (e.g., Mamba layers), to improve robustness under non-stationary artifacts and real-world noise, and to study whether end-to-end optimization further benefits open-set authentication.

% ============================================================
% REFERENCES
% ============================================================
\begin{thebibliography}{00}

\bibitem{b1} M. Fallahi, T. Strufe, and P. Arias-Cabarcos, ``BrainNet: Improving brainwave-based biometric recognition with Siamese networks,'' in \textit{Proc. IEEE PerCom}, 2023, pp. 53--60.

\bibitem{b2} H. Yang, W. Zheng, Y. Zhang, and B. L. Lu, ``EEG authentication based on deep learning of triplet loss,'' \textit{Neural Network World}, vol. 32, no. 5, pp. 269--283, 2022.

\bibitem{b3} X. Zhang, L. Yao, S. S. Kanhere, Y. Liu, T. Gu, and K. Chen, ``MindID: Person identification from brain waves through attention-based recurrent neural network,'' arXiv:1711.06149, 2017.

\bibitem{b4} N. A. Alzahab \textit{et al.}, ``Auditory evoked potential electroencephalography-biometric dataset,'' \textit{Data in Brief}, vol. 57, p. 111065, 2024.

\bibitem{b5} R. Hadsell, S. Chopra, and Y. LeCun, ``Dimensionality reduction by learning an invariant mapping,'' in \textit{Proc. IEEE CVPR}, 2006, pp. 1735--1742.

\bibitem{b6} F. Schroff, D. Kalenichenko, and J. Philbin, ``FaceNet: A unified embedding for face recognition and clustering,'' in \textit{Proc. IEEE CVPR}, 2015, pp. 815--823.

\bibitem{b7} J. Deng, J. Guo, N. Xue, and S. Zafeiriou, ``ArcFace: Additive angular margin loss for deep face recognition,'' in \textit{Proc. IEEE/CVF CVPR}, 2019, pp. 4690--4699.

\bibitem{b8} X. Wang \textit{et al.}, ``Multi-similarity loss with general pair weighting for deep metric learning,'' in \textit{Proc. IEEE/CVF CVPR}, 2019, pp. 5022--5030.

\bibitem{b9} S. Hochreiter and J. Schmidhuber, ``Long short-term memory,'' \textit{Neural Computation}, vol. 9, no. 8, pp. 1735--1780, 1997.

\bibitem{b10} A. Graves and J. Schmidhuber, ``Framewise phoneme classification with bidirectional LSTM,'' \textit{Neural Networks}, vol. 18, no. 5--6, pp. 602--610, 2005.

\bibitem{b11} K. He, X. Zhang, S. Ren, and J. Sun, ``Deep residual learning for image recognition,'' in \textit{Proc. IEEE CVPR}, 2016, pp. 770--778.

\bibitem{b12} D. L. Davies and D. W. Bouldin, ``A cluster separation measure,'' \textit{IEEE Trans. PAMI}, vol. 1, no. 2, pp. 224--227, 1979.

\bibitem{b13} Y. Luo and N. Mesgarani, ``TasNet: Time-domain audio separation network for real-time, single-channel speech separation,'' in \textit{Proc. IEEE ICASSP}, 2018, pp. 696--700, doi: 10.1109/ICASSP.2018.8462116.

\bibitem{b14} A. van den Oord \textit{et al.}, ``WaveNet: A generative model for raw audio,'' arXiv:1609.03499, 2016.

\end{thebibliography}

\end{document}
